% =============================================================
%  Résumé — Français, Anglais, Arabe (placeholder)
% =============================================================

% ── Résumé en français ────────────────────────────────────────
\chapter*{Résumé}
\markboth{Résumé}{Résumé}

\noindent\textbf{Titre~:} Conception et Déploiement d'une Application de Test
Telnet des Appareils Réseau avec Pipeline CI/CD DevOps.\\[0.2cm]
\noindent\textbf{Étudiant~:} Adem Hmercha \hfill
\textbf{Entreprise~:} Sagemcom Ezzahra\\[0.3cm]
\noindent\rule{\textwidth}{0.4pt}
\vspace{0.3cm}

\noindent Dans le cadre de la validation fonctionnelle des équipements réseau
embarqués développés par \textbf{Sagemcom Ezzahra} (passerelles résidentielles,
décodeurs multimédias, routeurs), les équipes de tests effectuent quotidiennement
des dizaines de commandes Telnet manuelles pour vérifier le comportement des
firmwares Linux. Cette approche artisanale génère des inefficacités opérationnelles
majeures~: absence de traçabilité, impossibilité d'automatiser les tests en parallèle
et manque de supervision centralisée.

Ce projet propose \textbf{Telnet Test Manager}, une application web fullstack
conçue pour centraliser, automatiser et historiser les tests Telnet des appareils
réseau. L'application est développée selon les principes de la \textit{Clean
Architecture} avec un backend \textbf{Node.js/Express.js}, une base de données
\textbf{MongoDB}, et un frontend \textbf{React~18/TypeScript}. Le moteur de tests
exploite les \textit{Worker Threads} de Node.js pour exécuter les sessions Telnet
en isolation, tandis qu'un serveur \textbf{WebSocket} diffuse les résultats en
temps réel. Sur le plan DevOps, l'application est entièrement conteneurisée avec
\textbf{Docker}, orchestrée via \textbf{Kubernetes} (chart Helm v0.3.0) avec
autoscaling horizontal, et déployée par un pipeline \textbf{GitHub Actions}
auto-hébergé. Une stack de supervision \textbf{Prometheus/Grafana} assure la
visibilité en temps réel des performances applicatives.

\vspace{0.3cm}
\noindent\textbf{Mots-clés~:} Telnet, Node.js, React, MongoDB, Docker,
Kubernetes, Helm, GitHub Actions, CI/CD, WebSocket, Clean Architecture,
Prometheus, Grafana, DevOps.

\vspace{0.8cm}
\noindent\rule{\textwidth}{0.4pt}

\vspace{0.8cm}

% ── Abstract in English ───────────────────────────────────────
\chapter*{Abstract}
\markboth{Abstract}{Abstract}

\noindent\textbf{Title:} Design and Deployment of a Telnet Testing Application
for Network Devices with a DevOps CI/CD Pipeline.\\[0.2cm]
\noindent\textbf{Student:} Adem Hmercha \hfill
\textbf{Company:} Sagemcom Ezzahra\\[0.3cm]
\noindent\rule{\textwidth}{0.4pt}
\vspace{0.3cm}

\noindent In the context of functional validation of embedded network equipment
developed by \textbf{Sagemcom Ezzahra} (residential gateways, set-top boxes,
routers), test teams manually execute dozens of Telnet commands daily to verify
Linux firmware behavior. This manual approach leads to significant operational
inefficiencies~: lack of traceability, inability to parallelize tests, and
absence of centralized monitoring.

This project introduces \textbf{Telnet Test Manager}, a fullstack web application
designed to centralize, automate, and archive Telnet tests for network devices.
Built on \textit{Clean Architecture} principles, the application features a
\textbf{Node.js/Express.js} backend, a \textbf{MongoDB} database, and a
\textbf{React~18/TypeScript} frontend. The Telnet engine leverages Node.js
\textit{Worker Threads} for isolated test execution, while a \textbf{WebSocket}
server streams real-time results. On the DevOps side, the application is fully
containerized with \textbf{Docker}, orchestrated via \textbf{Kubernetes}
(Helm chart v0.3.0) with horizontal autoscaling, and deployed through a
self-hosted \textbf{GitHub Actions} pipeline. A \textbf{Prometheus/Grafana}
monitoring stack provides live visibility into application performance.

\vspace{0.3cm}
\noindent\textbf{Keywords:} Telnet, Node.js, React, MongoDB, Docker, Kubernetes,
Helm, GitHub Actions, CI/CD, WebSocket, Clean Architecture, Prometheus, Grafana,
DevOps.

\vspace{0.8cm}
\noindent\rule{\textwidth}{0.4pt}
\vspace{0.5cm}

% ── Résumé en arabe (placeholder — nécessite XeLaTeX + arabxetex) ──
%
% Pour inclure un résumé en arabe, compilez avec XeLaTeX et remplacez
% cette section par le texte arabe en utilisant le package arabxetex
% ou polyglossia + une police arabique (Amiri, Scheherazade, etc.).
%
% \begin{RLtext}
%   ملخص: [texte arabe ici]
% \end{RLtext}

\begin{infobox}[title={Résumé en arabe}]
  \textit{Le résumé en arabe nécessite une compilation avec \textbf{XeLaTeX}
  et le paquet \textbf{polyglossia} avec une police arabique (Amiri ou
  Scheherazade). Remplacez cette boîte par le texte arabe correspondant
  une fois votre environnement XeLaTeX configuré.}
\end{infobox}
